\section{Transition surfacique--volumétrique et échelle spectrale}
\label{sec:ym-areal-volumetric-scale}

La constante qui fixe le semi-groupe de jauge est obtenue avant l'introduction de l'hamiltonien. Sa provenance combine deux lecteurs indépendants du même état généalogique : la fréquence de retour exprimant \(\pi_{\rm HMT}\) et la distribution stationnaire des deux feuillets aire--volume.

\subsection{Dynamique des deux feuillets}

Soit
\begin{equation}
 F_{\rm av}=\begin{pmatrix}0&1\\1&1\end{pmatrix},
 \qquad e_{\rm ar}=\binom10,
 \qquad e_{\rm vol}=\binom01.
 \label{eq:ym-av-transition-matrix}
\end{equation}
La récurrence de Fibonacci donne, par récurrence,
\begin{equation}
 F_{\rm av}^{n}
 =\begin{pmatrix}F_{n-1}&F_n\\F_n&F_{n+1}\end{pmatrix}.
 \label{eq:ym-av-fibonacci-powers}
\end{equation}
Le vecteur de Perron est proportionnel à \((1,\varphi)\). Comme
\(F_{\rm av}\) est symétrique, la mesure de Parry assigne aux deux feuillets
des poids proportionnels à \((1,\varphi^2)\), et donc
\begin{equation}
 \frac{\mu_{\rm ar}}{\mu_{\rm vol}}=\varphi^{-2},
 \qquad
 \mu_{\rm vol}=\frac{\varphi^2}{1+\varphi^2}.
 \label{eq:ym-av-parry-weights}
\end{equation}
Ici, \(\mu_{\rm vol}\) est un caractère de fréquence sans dimension ; il ne désigne pas un volume de dimension \(L^4\).

Le lecteur régional de \(\pi\) est introduit une seule fois au moyen de
\begin{equation}
 D_\pi
 =\pi_{\rm HMT}|e_{\rm ar}\rangle\langle e_{\rm ar}|
  +|e_{\rm vol}\rangle\langle e_{\rm vol}|.
 \label{eq:ym-pi-regional-reader}
\end{equation}
Le rapport exprimé à la profondeur \(n\) est
\begin{equation}
 \Theta_n
 :=\frac{\langle e_{\rm ar},D_\pi F_{\rm av}^{n}e_{\rm ar}\rangle}
 {\langle e_{\rm vol},F_{\rm av}^{n}e_{\rm vol}\rangle}
 =\pi_{\rm HMT}\frac{F_{n-1}}{F_{n+1}}
 \longrightarrow
 r_{\rm av}:=\frac{\pi_{\rm HMT}}{\varphi_{\rm HMT}^{2}}.
 \label{eq:ym-target-free-av-limit}
\end{equation}
La matrice, le chemin et le quotient sont fixés avant de consulter un saut de masse ou un paramètre de théorie de jauge. À la profondeur native \(n=108\),
\begin{equation}
 F_{107}=10284720757613717413913,
 \qquad
 F_{109}=26925748508234281076009,
 \label{eq:ym-fibonacci-107-109}
\end{equation}
et
\begin{equation}
 \left|\Theta_{108}
 -\frac{\pi_{\rm HMT}}{\varphi_{\rm HMT}^{2}}\right|
 <2\cdot10^{-45}.
 \label{eq:ym-av-depth-108-error}
\end{equation}

\subsection{Transduction dimensionnelle}

Soit \(\ell_0\) l’unité longitudinale du lecteur dimensionnel. L’échelle
inverse de longueur et le paramètre d’une dalle d’épaisseur \(a_m\) sont
\begin{equation}
 \kappa_{\rm av}
 :=\frac{\mu_{\rm vol}}{\ell_0}
 \left(\frac{\pi_{\rm HMT}}{\varphi_{\rm HMT}^{2}}-1\right)>0,
 \qquad
 q_m:=e^{-3a_m\kappa_{\rm av}}.
 \label{eq:ym-dimensional-gap-scale}
\end{equation}
Par conséquent, \([\kappa_{\rm av}]=L^{-1}\) et l’écart énergétique est
\begin{equation}
 \Delta_E=3\hbar c\,\kappa_{\rm av};
 \label{eq:ym-energy-gap-units}
\end{equation}
en unités \(\hbar=c=1\), on écrit \(\Delta_E=3\kappa_{\rm av}\).
La mémoire électronique \(\Delta_4\) appartient à un autre domaine et à un autre
lecteur ; elle n’intervient pas dans \eqref{eq:ym-dimensional-gap-scale}.

Dans les semi-groupes des sections suivantes, \(t,s,a_m\) sont des longueurs euclidiennes et les générateurs \(D\) et \(H^{\rm ov}\) ont pour dimension \(L^{-1}\). L'hamiltonien énergétique est \(H_E=\hbar_{\rm int}c_{\rm int}D\). Si \(\tau=t/c_{\rm int}\), alors
\[
 e^{-tD}=e^{-\tau H_E/\hbar_{\rm int}},\qquad
 \Delta=3\kappa_{\rm av},\quad
 \Delta_E=\hbar_{\rm int}c_{\rm int}\Delta,\quad
 m_{\rm gap}=\frac{\hbar_{\rm int}\Delta}{c_{\rm int}}.
\]
Le même saut est ainsi exprimé respectivement comme longueur inverse, énergie et masse dans la coordinatisation dimensionnelle déclarée.

Comme \(a_m=9a_{m+1}\), la définition donne immédiatement
\begin{equation}
 q_{m+1}^{9}=q_m.
 \label{eq:ym-nonadic-gap-root}
\end{equation}
Cette identité est la racine nonadique qui transportera le même semi-groupe et la
même échelle spectrale à travers tous les raffinements.
